\chapter{Perito}
\label{cap:perito}

O perito selecionado para apoiar este projeto foi o Professor Doutor Agostinho
Gil Teixeira Lopes, Professor Adjunto do Departamento de Engenharia Informática
do Instituto Superior de Engenharia do Porto (ISEP). A sua formação e
experiência abrangem áreas diretamente relacionadas com o projeto, nomeadamente
robótica, sistemas autónomos, visão por computador, inteligência artificial,
sensores, automação e controlo \preencher{fontes da biografia}.

Concluiu a licenciatura em Informática e Matemáticas Aplicadas em 2002, na
Universidade Portucalense, e o doutoramento em Engenharia Informática em 2006,
na Cranfield University, no Reino Unido, com especialização em automação,
controlo e robótica. Desenvolveu posteriormente atividades de investigação e
pós-doutoramento relacionadas com a automatização e o controlo de processos
industriais \preencher{fonte}.

Ao longo da sua carreira, exerceu funções de investigação e docência na
Universidade do Minho e na Universidade da Maia, integrando o ISEP em outubro de
2024. A sua atividade científica inclui investigação em robótica móvel autónoma,
visão artificial, redes neuronais, localização e navegação de robôs, bem como
orientação de projetos académicos nestas áreas \preencher{fonte}.

A sua experiência justifica a escolha como perito, uma vez que lhe permite
avaliar a representação e interpretação dos dados sensoriais, os critérios de
criticidade e as regras de tomada de decisão do sistema. O seu contributo é
particularmente relevante para validar a lógica de decisão e os mecanismos de
explicação implementados no projeto.
